% ---------- Classification (Krempl part 4) ----------
\begin{topic}{8 \; Prediction: classification}{cClass}
\insight{Task: learn classifier $f:X\to Y$ from labelled training sample (features $X$, class label $Y$) that optimises a performance measure on \emph{new} data, e.g.\ 0-1 loss $L(y,f(x))=0$ if $f(x)=y$, else 1. Probabilistic classifiers also return posterior $P(y|x)$.}
\sub{Probability rules}
\formula{Joint}{P(X,Y)=P(X|Y)P(Y)=P(Y|X)P(X)}{}
\formula{Marginalise}{P(x)=\sum_{y}P(x,y)}{sum over the values of \term{y} (not $x$)! Probabilities $\ge0$, sum to 1; mutually exclusive events add.}
\formula{Bayes theorem}{P(Y|X)=\frac{P(X|Y)\,P(Y)}{P(X)}}{posterior = class-conditional $\times$ prior / evidence.}
\formula{Prior}{P(Y{=}+)=\frac{|L_+|}{|L_+\cup L_-|}}{}
\insight{\term{Bayes-optimal}: predict $+$ iff $P(+|x)>P(-|x)$ $\Leftrightarrow$ $P(x,+)>P(x,-)$. Decision boundary where they are equal (intersections of the joint curves). \term{Bayes error rate} = unavoidable error: at each $x$ the smaller posterior $\min_y P(y|x)$; integrate (area under the lower joint curve). Zero FPs \emph{and} zero FNs impossible when classes overlap.}
\ex{$P(x{=}4,+)=.077$, $P(x{=}4,-)=.159$ $\to P(x{=}4)=.236$, $P(+|4)=.077/.236=.33<.67\to$ classify $-$. In a region where everything is called $+$: FN rate $=0$.}
\sub{Estimating $P(X|Y)$}
\insight{\term{Histogram}: discretise $X$ into intervals, count per class, frequencies as $P(x|y)$. Interval too big $\to$ loss of detail; too small $\to$ lack of data (empty cells).}
\formulaD{Kernel density estimate}{K(x,x_i)=\frac{1}{\sqrt{2\pi\sigma^2}}\exp\Big(-\frac{(x-x_i)^2}{2\sigma^2}\Big),\qquad \hat P(x|+)=\frac{\sum_{x_i\in L_+}K(x,x_i)}{|L_+|}}{Gaussian on every training instance; bandwidth $\sigma$ small $\to$ overfits noise/outliers, large $\to$ oversmooths. Repeat per class, multiply by prior, compare.}
\sub{Multivariate data}
\formula{Multivariate Bayes}{P(Y|X_1,X_2)=\frac{P(X_1,X_2|Y)P(Y)}{P(X_1,X_2)}}{needs the \emph{joint} $P(X_1,X_2|Y)$ (counts, or parametric model e.g.\ Gaussian mixture) $\to$ curse of dimensionality ($2^D$ cells).}
\formula{Naive Bayes}{P(Y|X_1..X_p)\propto P(Y)\prod_{i=1}^p P(X_i|Y)}{assumes \term{conditional independence} of features given the class: $P(X_i|X_j,Y)=P(X_i|Y)$. Violated when, within a class, $X_1$ and $X_2$ co-vary (check $P(X_1,X_2|Y)$ vs $P(X_1|Y)P(X_2|Y)$ per cell).}
\ex{Golf (14 rows, 9 yes): $P(\text{hot})=4/14$, $P(\text{yes})=9/14$, $P(\text{hot},\text{yes})=2/14$, $P(\text{hot}|\text{yes})=2/9$, $P(\text{yes}|\text{hot})=2/4$. Cool \& normal: multivariate $=3/4$ yes. Naive: yes $\propto\frac9{14}\cdot\frac39\cdot\frac69=\frac17$, no $\propto\frac5{14}\cdot\frac15\cdot\frac15=\frac1{70}$ $\to$ yes (normalised $10/11$).}
\sub{Other classifiers}
\insight{\term{k-NN}: memorise all training data; for new $x$ compute distance to all, assign majority label of $k$ nearest. \term{Lazy learner}: no model/abstraction, fast training, slow prediction.}
\insight{\term{Decision tree}: iteratively split on features, greedy choice of split attribute (e.g.\ information gain); leaf predicts its majority class.}
\insight{\term{Generative} (models the data distribution, can generate data, e.g.\ Bayes) vs \term{predictive} (learns the decision rule directly; simpler problem). \term{Ensembles} (multiple classifiers): more robust \& accurate, but more time/space, less interpretable.}
\insight{All classifiers make assumptions: naive Bayes (cond.\ independence), clustering assumption, linearity (regression-based).}
\sub{Experimental design}
\insight{\term{Training set}: fit classifier. \term{Validation set}: compare parameter values / models. \term{Test set}: final, realistic out-of-sample estimate. Never evaluate on training data: \term{overfitting} = model fits noise/outliers, generalises badly.}
\insight{\term{k-fold cross-validation}: split into $k$ subsets; each is test set once, the others train; good for small data, computationally expensive.}
\sub{Performance measures}
\insight{Confusion matrix (slides): \term{columns = actual}, \term{rows = predicted}; diagonal correct. \term{FP}: actually $-$, predicted $+$. \term{FN}: actually $+$, predicted $-$. Can be weighted by a cost matrix.}
\formula{Accuracy}{\frac{TP+TN}{TP+TN+FP+FN}}{Error rate $=\frac{FP+FN}{\text{all}}=1-\text{Acc}$; misleading under class imbalance.}
\formula{Sensitivity}{\text{TPR}=\frac{TP}{TP+FN}}{share of actual positives found.}
\formula{Specificity}{\frac{TN}{TN+FP}=1-\text{FPR},\ \ \text{FPR}=\frac{FP}{FP+TN}}{share of actual negatives correctly rejected.}
\formula{Misclass.\ loss}{FP\cdot\text{Cost}_{FP}+FN\cdot\text{Cost}_{FN}}{for unequal error costs; often $\text{Cost}_{FP}+\text{Cost}_{FN}=1$ (= cost ratio).}
\insight{\term{ROC}: TPR (y) vs FPR (x) over all thresholds; steeper rise = better. \term{AUC} = area: 1 perfect, 0.5 diagonal = random guessing. A curve that is everywhere above another \term{dominates} it; crossing curves $\Rightarrow$ neither dominates, best one depends on FP/FN trade-off. Suited for class imbalance.}
\sub{Caveats}
\insight{Every single measure implicitly assumes an FP/FN trade-off $\Rightarrow$ ask for the cost ratio. Imbalance: cost-sensitive classifier or re-weight/re-sample \emph{before} training (oversample minority, undersample majority); correcting predictions afterwards is a bad idea.}
\insight{Past as model for the future: assumes training data representative; drift $\Rightarrow$ monitor \& update. \term{Observer paradigm}: model assumed not to influence environment; feedback loops break it (spammers adapt, bad credit score raises default risk).}
\insight{Labels scarce, features cheap: \term{semi-supervised} (use unlabelled data + structure assumptions), \term{active learning} (choose which instances to label).}
\end{topic}
